\DTMsavedate{mydate}{2022-06-30}
\maketitle{Integrated Circuits and Microelectronics}{Electronic Engineering}{\DTMusedate{mydate}}

% ----------------------------------------------------- %
% COURSE INFO
\subsection*{Course Information}\label{course:5111}
\vspace{0ex}%
\noindent\parbox{0.5\textwidth}{%
    \noindent\begin{tabular}{@{}ll}
                 \textsf{Course:}          & 	Integrated Circuits and Microelectronics  \\
                 \textsf{Course Code:}         & 	TEE 5111    \\
                 \textsf{Credits:}    & 	10
    \end{tabular}}
\vspace{2ex}\hrule\vspace{2ex}

% ----------------------------------------------------- %
% INSTRUCTOR INFO
\subsection*{Lecturer Information}
\noindent\parbox{0.5\textwidth}{%
    \noindent\begin{tabular}{@{}ll}
                 \textsf{Name:}         &    Svetlana A. Bebova           \\
% \textsf{Phone:}        &    \phone          \\
\textsf{Email:}        &    \href{mailto:swetlana.bebova@nust.ac.zw}{swetlana.bebova@nust.ac.zw} \\
\textsf{Qualifications:}        &  MSc in Radio Electronic Engineering (Technical University of Sofia)
    \end{tabular}}
\vspace{2ex}\hrule\vspace{2ex}

% ----------------------------------------------------- %
% COURSE SYNOPSIS
\subsection*{Course Synopsis}
The module gives a comprehensive view of operational amplifiers i.e.\ classification, parameters, basic building blocks and data sheets, thermal analysis, applications.
It also delves into differential, instrumentation and isolation Amplifiers i.e.\ its parameters, analysis, data sheets and applications.
It also looks into data converters such as ADC, DAC, V/F and F/V converters i.e.\ its operation, analysis and applications.
Elements of data acquisition systems such power supply systems, over-voltage and short circuit protection are covered in this course including microelectronic technologies as well.
\vspace{2ex} \hrule\vspace{2ex}

% ----------------------------------------------------- %
% LEARNING OBJECTIVES
% \subsection*{Learning Objectives}
% By the end of the course, candidates should be able to:
% \begin{outline}
%    \1      Understand the concept of states as applied to Control Systems.
% \end{outline}
% \vspace{2ex}\hrule\vspace{2ex}

% ----------------------------------------------------- %
% COURSE TOPICS
\subsection*{Topics}
\begin{outline}
    \1 Operational amplifiers and linear integrated circuits (Ics) classification and parameters.
    \1 Ics basic building blocks
    \1 Microelectronics:
    \begin{itemize}
        \item Wafer preparation
        \item Impurity doping
        \item Oxidation
        \item Photolithography
        \item Metalization
    \end{itemize}
    \1 Integrated components
    \1 Special types amplifiers
    \begin{itemize}
        \item Differential amplifier (integrated and monolithic)
        \item Instrumentation amplifier
        \item Isolation amplifier
    \end{itemize}
    \1 Thermal analysis
    \1 Introduction to FPGA
    \1 Data converters
    \begin{itemize}
        \item Digital-to Analog
        \item Analog–to Digital
        \item Voltage-to-frequency and Frequency-to-Voltage
    \end{itemize}
\end{outline}
\vspace{2ex} \hrule\vspace{2ex}

% ----------------------------------------------------- %
% Students Assenssment
\subsection*{Assessment of Students}
\begin{outline}
    \1 Continuous assessment will consist of tests and assignments.
    \1 Final assessment will consist of an examination at the end of the semester.
\end{outline}
\vspace{2ex}\hrule\vspace{2ex}

% ----------------------------------------------------- %
% Grade Evaluation
\subsection*{Course Evaluation and Grading Policies}
Grades are based on:
Continuous assessment (\qty{25}{\percent}),
Final Exam (\qty{75}{\percent})
% \vspace{2ex}\hrule\vspace{2ex}
% ----------------------------------------------------- %
% EQUIPMENT
% https://sites.udel.edu/computing-purchases/personal-specs/
% \subsection*{Essential Equipment}

% \vspace{2ex}\hrule\vspace{2ex}

% Policies and Other information
% \subsection*{Policies and Other Information}
% Key Text or some Books
% \subsection*{Key Text}

% \begin{itemize}
%   \item Roland S Burns, Advanced Control Engineering
%   \item D K Anand, R B Zmood, Introduction to Control Systems
%   \item E A Parr, Control Engineering
% \end{itemize}
\vspace{2ex}\hrule
\vspace{2ex}\hrule\vspace{2ex}